\section*{Basics}

\subsection*{Mathematik}
\subsubsection*{Rechnen mit Wahrscheinlichkeiten}
$P(A \cup B) = P(A) + P(B)$, wenn $A,B$ disjunkt.

$P(A \text{ und } B) = P(A) \cdot P(B)$, wenn $A,B$ unabhängig.


\subsubsection*{Bedingte Wahrscheinlichkeit}
$\displaystyle P(X \mid C) = \frac{P(X \cap C)}{P(C)}$

\subsubsection*{Unabhängigkeit}
$P(X \cap Y) = P(X) \cdot P(Y)$

\subsubsection*{Erwartungswert und Varianz}
$\displaystyle \mathbb{E}(X) = \sum_i x_i \cdot P(X = x_i)$

$\displaystyle \text{Var}(X) = \mathbb{E}[(X-\mathbb{E}(X))^2] = \mathbb{E}(X^2) - \mathbb{E}(X)^2$

\subsection*{Modelle}
Hypothesenklasse: Menge an Funktionen

\subsubsection*{Parametrisierte Modelle}
Parameter $\theta$ definiert Funktion im Hypothesenraum

Finde $\theta$ mit minimalem Loss auf Daten.

\subsubsection*{Beispiel: Rechtecksklassifikation}
Hypothesenklasse $\mathcal{H}$: Menge aller Rechtecke.

$h\in\mathcal{H}$ definiert durch $\theta=(o_1,o_2,u_1,u_2)$.

\subsubsection*{Generalisierung}
Wie gut ist ein Modell auf neuen Daten.

Underfitting: Komplexität $\mathcal{H} <$ Daten

Overfitting: Komplexität $\mathcal{H} >$ Daten

\subsubsection*{No-Free-Lunch Theorem}
Ein Algorithmus kann nur für ein spezifisches Problem geeignet sein.

Es gibt nicht den besten Algorithmus.

\subsection*{Modellselektion}
Vermeide: Ganzen Datensatz benutzen (Memorisierung)

\subsubsection*{Test- und Validierungsset}
Teile das Dataset auf in:\\
Trainingsset, Validierungsset, Testset

\subsubsection*{$K$-fold Kreuzvalidierung}
Aufteilung: $K$x Training/Validierung, 1x Test.

Nutze nacheinander eins der $K$ Sets als Validation, Mittelung über alle Durchführungen.
